\chapter{Plantar Fasciitis}
\index{Plantar Fasciitis}

\medskip
\noindent\textit{This chapter is for general education. Persistent foot pain should be assessed, particularly when diabetes or reduced sensation is present.}

\section*{What it is}
Plantar fasciitis is pain associated with irritation or overload of the plantar fascia, a strong band running along the sole of the foot. Pain is often felt near the heel or arch and may be worse with the first steps after sleep or rest.

\section*{Contributing factors}
Risk can increase after a sudden change in activity, prolonged standing or walking, tight calf muscles, unsupportive footwear, running on hard surfaces, or higher body weight. Similar pain can have other causes, so a symptom pattern alone does not confirm the diagnosis.

\section*{Self-care and treatment}
Reduce aggravating activity temporarily without complete inactivity, use supportive cushioned footwear, and try gentle calf and sole stretches. Ice wrapped in a cloth may help short-term pain. A pharmacist or clinician can advise on pain medicine, insoles, night splints, physiotherapy, or other treatment. Steroid injections and surgery are reserved for selected persistent cases.

\section*{When to Seek Medical Care}
Arrange assessment if pain is severe, worsening, recurrent, prevents normal activities, has not improved after a short period of self-care, or is associated with tingling or loss of feeling. People with diabetes should seek advice early for new foot pain or skin changes. Seek urgent care after major injury, if the foot is cold or pale, or if there is rapidly spreading redness or fever.

\section*{Sources and further reading}
\begin{itemize}
\item NHS. \textit{Plantar fasciitis}. \url{https://www.nhs.uk/conditions/plantar-fasciitis/}
\end{itemize}
